\begin{frame}{그림으로 읽기: 커널 하나만 바꿨을 뿐}
  \begin{columns}[T]
    \begin{column}{0.52\textwidth}
      \centering
      \includegraphics[width=\linewidth]{ch05_3_kernel_boundaries.png}
    \end{column}
    \begin{column}{0.44\textwidth}
      \small
      \begin{itemize}
        \item \kb{선형 커널}(좌): 동심원을 직선 하나로 어정쩡하게
          자를 뿐 — 정확도 \textbf{0.52}(동전 던지기 수준).
        \item \kb{RBF 커널}(우): \textbf{원형 결정 경계}로 안쪽·
          바깥쪽 원을 완벽히 분리 — 정확도 \textbf{1.00}.
      \end{itemize}
      \vskip0.4em
      \begin{block}{서포트 벡터 개수가 핵심}
        선형 커널: 학습 데이터 \textbf{137개}$(\approx68\%)$를
        SV로 끌고 옴.
        RBF 커널: \textbf{31개}만 SV로 쓰고 나머지 169개는
        마진 바깥에 둠.
        \vskip0.3em
        5.2절의 ``서포트 벡터가 경계를 떠받친다''는 원칙이
        \textbf{비선형 경계에도 그대로} 적용되는 직접적 증거.
      \end{block}
    \end{column}
  \end{columns}
\end{frame}
